\section{Introduction}

District-heating substations operate continuously and generate large volumes of time-series data, but real fault labels are often sparse or unavailable. This creates a practical detection problem: operators still need a way to identify unusual operating periods, even when a supervised classification setup is not possible. The aim of this thesis is therefore to build and evaluate an anomaly-detection pipeline that can identify physically interpretable abnormal behavior in district-heating substations using historical data and later be applied to live data.

The work is motivated by the HEAT paper, which combines encoder-based representation learning and clustering for fault detection in district-heating substations. However, the available KTH thesis dataset differs in two important ways. First, the network is much smaller, which weakens the statistical value of peer-group clustering. Second, the historical coverage is much longer, which makes per-substation temporal modeling more promising than network-wide peer comparison alone. Because of this, the thesis focuses primarily on reconstruction-based anomaly detection and treats clustering as a supporting interpretation tool rather than the central detection mechanism.

The main research question is:

\begin{quote}
Can reconstruction-based anomaly detection on supply temperature, return temperature, and derived flow identify interpretable abnormal operating periods in a small district-heating network?
\end{quote}

The thesis also addresses the following supporting questions:

\begin{itemize}
\item How should the derived flow feature be handled when small temperature differences cause numerical instability?
\item How does a reconstruction-based detector compare with a simple engineering low delta-T baseline?
\item What is the effect of threshold choice on the number and type of detected anomalies?
\item Can clustering still help as an interpretation layer when the number of substations is small?
\end{itemize}

The rest of the thesis is organized as follows. Chapter~2 reviews the relevant background and the HEAT paper. Chapter~3 describes the available datasets and preprocessing assumptions. Chapter~4 presents the anomaly-detection methodology. Chapter~5 summarizes the experimental results. Chapter~6 discusses interpretation, limitations, and implications for live deployment. Chapter~7 concludes the thesis and outlines future work.
